% Extensión estimada: una página, cuatro párrafos (ficha 5.6).
% Fuentes: secciones 5.1 a 5.5. Sin números nuevos.
% Feedback del tutor (14/09/2026): el cierre se llama "Discusión"; la discusión que cruza clasificación y generación
% queda en el Capítulo 6.
\section{Discusión de los resultados de generación}
\label{sec:res_gen_discusion}

Los resultados del capítulo permiten responder la segunda pregunta de investigación
(Sección~\ref{sec:intro_objetivos}) en sus dos partes. En cuanto al diagnóstico, un \gls{vlm} generativo ajustado no
afirma por sí solo la clase de la imagen: sin un diagnóstico, enumera clases, evita comprometerse o, en el caso del
generador especializado, repite un texto ajeno a las siete clases. Con la cascada, en cambio, el reporte afirma el
diagnóstico del clasificador, y su fidelidad diagnóstica coincide con la exactitud de Swin-Tiny, tanto en total como
en cada clase. Es decir, el reporte afirma el diagnóstico correcto en la medida en que el clasificador lo acierta, y
ningún generador corrige un diagnóstico equivocado mirando la imagen.

En cuanto a la descripción, la imagen cambia lo que el reporte describe de cada lesión, pero no se pudo verificar que
lo descrito esté en ella. Sin la imagen, los generadores escriben un único texto por clase, y con ella cada reporte es
distinto, con más entidades clínicas en el caso de Mini-VLM. Sin embargo, las métricas automáticas miden la presencia
de términos y no su validez, y el ejemplo del melanoma afirmado como nevo muestra que la descripción también se adapta
al diagnóstico recibido. Por lo tanto, la segunda parte de la pregunta queda respondida solo en lo que las métricas
sin especialista pueden establecer.

El ruteo por concordancia agrega al \emph{pipeline} lo que la cascada no tiene, una forma de reconocer los casos
dudosos. La concordancia entre Swin-Tiny y el \gls{knn} separa los casos en que el clasificador es confiable, y en los
demás presentar dos candidatos cubre más que afirmar uno. La mejora del reporte, sin embargo, es del ruteo y no del
generador: Mini-VLM la conserva y LLaVA-Med la pierde cuando razona más allá del límite de tokens. Además, los errores
que comparten los dos predictores, entre ellos casi un cuarto de los melanomas, se siguen afirmando como un acierto.

Entre los generadores, Mini-VLM resulta el más adecuado para el \emph{pipeline}: transcribe el diagnóstico como el
generador especializado, describe tanto como el biomédico más grande y es el único cuyos pesos entran en una GPU de
consumo de 12~GB. En conjunto, el diagnóstico del reporte proviene del clasificador, la descripción proviene del generador, y
el ruteo decide cuánto afirma el \emph{pipeline}. El Capítulo~\ref{cp:discusion} discute las limitaciones de esta
evaluación junto con las de la clasificación, en particular la ausencia de una revisión por especialistas y el uso de
una sola semilla para el clasificador que ancla los reportes.
